\section{Further research questions and proposed next steps}
\label{sec:research}

The following questions are ordered roughly from local implementation work
to the remaining global topology barriers. Each has a concrete success
criterion. They are proposed tasks, not consequences of the current proofs.

\subsection{Make least-transversal replay closer to linear in trace size}

Can a persistent ordered interval sequence replace the current repeated
live-list scans and reverse unions, with an
$O((E+G)\log(G+2))$ structural bound? The least-representative theorem and
$G$-interval output bound already provide the necessary compact output.
The missing work is a data structure supporting rank splits, gap insertion,
prefix truncation, and coalescing with explicit binary costs and a simple
checker. A successful implementation should show both a proof of the
improved structural bound and a crossover measurement against the small
list implementation; asymptotically better trees may lose on short traces.

\subsection{Compile reusable weighted observers}

The same normal-arc relation can receive Euler, boundary, coordinate, or
other additive weights. The current weighted kernel replays a validated
orbit trace for each query. Can that trace be compiled into a compact
operator that transports new interval weights without repeating unrelated
arithmetic or source reconstruction? The operator must retain the
contraction/folding structure and avoid a dense matrix over the expanded
point universe. A useful milestone is a source-bound transport circuit
with a polynomial construction bound and a query bound in its explicit
size, measured on sequences of distinct observers of the same surface.

\subsection{Add essential-disc information without losing correlations}

On a torus boundary, the maintained two integral boundary-intersection
weights distinguish a compressing disc from a vertex-link disc. Combining
them with $(\chi,b)$ yields a four-coordinate integer observer. The
cover inversion extends verbatim to $\Z^4$, after which boundary fields
can be reduced modulo two. Can a verified four-weight type census replace
separate topology and essentiality calls efficiently? The test must retain
the correlation of boundary class with a particular disc type; combining
unlabelled totals afterward is unsound. Any attempt to use torsion weights
during inversion should supply new preimage-sum equations and prove their
uniqueness or describe the remaining ambiguity.

\subsection{Retain the representative-to-type association}

The present histogram forgets which least representative carries which
type. Weighted contraction events and the live-source interval map could
instead emit a piecewise-constant signature function on the transversal.
Can that association be retained with a polynomial number of intervals,
and then used to request the normal coordinates of one selected component?
This would connect an inexpensive topology census to a concrete surface
needed for cutting. The required output is a selected vector with a
source-bound component witness, not an arbitrary vector having the same
Euler characteristic and genus.

\subsection{Extend the geometry validator beyond one torus boundary}

The two-weight construction itself does not need a boundary homology basis.
Its next natural domain is compact orientable manifolds with several
boundary components of arbitrary genus, as arise after cutting in a
hierarchy. The immediate goal is an independently audited validator for
that domain, with normal boundary-arc inclusion and the same source-bound
certificates. This generalization would not by itself recognize
compressing discs: on a higher-genus boundary, an essential separating
curve has zero homology class. A correct essentiality stage needs compressed
complement topology or a suitable normal-curve algorithm, rather than an
unchanged torus parity test. The fast curve algorithms of Lackenby provide
relevant primitives and quantitative models \cite{lackenby-curves}.

\subsection{Extract boundary slope or attachment data in compressed form}

Knowing that a component has $b$ boundary circles is weaker than knowing
their slopes and how they meet a boundary pattern. On a torus, can one
compute a compressed distribution of oriented primitive slopes, associated
with component representatives, without listing parallel circles? For
general boundary patterns, the corresponding target is ordered attachment
data with a controlled encoding. Both queries require more structure than
the scalar marker count. A successful interface should state exactly what
gluing operations it supports and certify the association to the original
surface.

\subsection{Find larger geometric cores than a single common factor}

The current core removes all vertex links and one global quadrilateral
content. Many surfaces contain several large parallel families with
different multiplicities and total coordinate gcd one. Can these be
decomposed into a bounded number of independently certified interval-bundle
families plus a small residual? A coordinate sum alone is not enough:
normal addition can change connectivity, so each independently scaled
family needs a geometric disjointness or bundle-monodromy certificate.
The repository's earlier dihedral-cover methods offer a restricted
starting point. A promising theorem would bound the residual encoded
complexity by the number of genuinely interacting families.

\subsection{Complete certified diagram-to-exterior provenance}

The normal-surface validator certifies a triangulated manifold and a vector
in it. To use a disc certificate for a knot diagram, a separate construction
must establish that this is its exterior, with the correct boundary data.
The next milestone is a reproducible diagram-to-triangulation producer and
a small combinatorial provenance checker whose moves preserve the specified
complement. Practical topology systems can assist construction, but an
unverified call returning an isomorphism signature is not the missing
proof. The certificate should survive subsequent triangulation simplification.

\subsection{Discover useful surfaces with a controlled search parameter}

The central unresolved search question remains: which normal vectors should
be supplied? Could bounded frontier width, a small family of quadrilateral
supports, or a compressed hierarchy invariant guarantee a useful disc or
multi-surface among only quasi-polynomially many candidates? A useful
restricted theorem must specify the input class, give an effective
candidate procedure, and charge arithmetic in coordinate bits. Merely
assuming logarithmic width after seeing successful examples would not
establish a general bound. The present census can certify the topology
of proposed candidates, making it suitable for testing such search rules.

\subsection{Implement a compressed cut with provenance}

An important next geometric operation is cutting along a selected normal
surface while retaining a polynomial description of the resulting cells,
boundary pieces, and attachment maps. The topology spectrum can decide
which component types are present, but it does not specify this cut.
A defensible first target is a restricted interval-bundle or layered
normal family where the attachment map has an explicit small grammar.
The implementation should produce a source-bound gluing certificate and
a theorem bounding the size of every intermediate representation. This
would advance the hierarchy programme more directly than another
uncorrelated numerical invariant.

\subsection{Control hierarchy depth, resets, and total encoded volume}

Once local cuts and surface discovery are available, the remaining task
is an amortized global argument. Can one define an explicit potential
whose decrease or controlled growth bounds the number of compressions,
discarded later surfaces, and hierarchy restarts? The potential must also
control encoded attachment data, not just Euler characteristic or the
number of handles. A concrete outcome would specify the parameters of
Lackenby's accelerated multi-surface/Cheeger-region programme well enough
to prove a product bound on all local operations. Distinguish a bound
$2^{O((\log n)^3)}$ from the stronger
$2^{O((\log n)^2)}$ goal when reporting progress.

\subsection{Formalize the finite combinatorial certificates first}

The least-transversal invariant, reverse gap formula, and finite-support
cover inversion are suitable initial Lean formalization targets. Their
objects are finite relations, interval maps, and integer histograms;
they can be stated without formalizing all of three-manifold topology.
A staged project could first verify the algebraic equations and the
selector compiler, then connect them to formal normal-arc and orientation
cover constructions. The delivered Python tests and certificates are
valuable evidence but are not machine-checked mathematical proofs.

\subsection{Measure an integrated geometric workload}

The present timings concern supplied-vector topology, not elapsed time to
recognize a knot from a diagram. Once a discovery stage actually invokes
these observers, measure the number of calls, the amount of core reduction,
and the total fraction of runtime they consume on a held-out corpus.
An optional observer portfolio could choose direct versus reduced versus
coordinate-rich queries from measured structural features. Any policy
should preserve exact certificates and publish its negative cases. The
success criterion is a faster complete workload with the same mathematical
decision, rather than a faster isolated function that rarely occurs.

\Needspace{10\baselineskip}
\section{Concluding assessment}

The implemented result is a complete compressed topology spectrum for a
supplied normal surface, with a constant-dimensional observer and an
orientation-aware numerical core. It includes a certified least-transversal
primitive, actual boundary-circle counts, sparse cover inversion, and
independent original-source replay. These operations remove identifiable
representation and repeated-arithmetic costs and give useful building
blocks for later geometric processing.

Surface discovery, compressed cutting with attachment data, and global
hierarchy accounting remain the decisive next steps toward quasi-polynomial
recognition. The coordinate route remains available when an embedded
component is required.
